% ============================================================================
% WORKBOOK — Additional: Proportional Reasoning with Scale Models
% State-specific (OK)
% Practice-focused reinforcement for the study guide topic
% ============================================================================
\section{Proportional Reasoning with Scale Models}
\topicTitle{Proportional Reasoning with Scale Models}

\begin{quickReview}{Scale Models}
A \textbf{scale model} is a smaller (or larger) copy of a real object. The \textbf{scale factor} tells you how model measurements relate to real ones.

\begin{itemize}[leftmargin=*]
    \item Scale $1 : 50$ means every $1$ unit on the model equals $50$ units in real life.
    \item \textbf{Actual length} $=$ model length $\times$ scale factor
    \item \textbf{Model length} $=$ actual length $\div$ scale factor
\end{itemize}

\medskip
\textbf{Example:} A model car is $8$ cm long. Scale is $1 : 25$. Actual length $= 8 \times 25 = 200$ cm $= 2$ m.
\end{quickReview}

% ── Warm-Up ──────────────────────────────────────────────────────────────────
\begin{practiceBox}[funGreen]{\faSun~Warm-Up}

\practiceHeader[funGreen]{Model to Actual (Multiply by the Scale Factor)}

\begin{multicols}{2}
\prob A model is $5$ cm long. Scale: $1 : 100$. Find the actual length. \answerBlank[3cm]
\answerExplain{$500$ cm}{Multiply the model length by the scale factor: $5 \times 100 = 500$ cm. The actual object is $500$ cm ($5$ m) long.}

\prob A model is $3$ cm wide. Scale: $1 : 200$. Find the actual width. \answerBlank[3cm]
\answerExplain{$600$ cm}{Multiply: $3 \times 200 = 600$ cm. The actual width is $600$ cm ($6$ m).}

\prob A model tower is $10$ cm tall. Scale: $1 : 50$. Find the actual height. \answerBlank[3cm]
\answerExplain{$500$ cm}{Multiply: $10 \times 50 = 500$ cm. The actual tower is $500$ cm ($5$ m) tall.}

\prob A model boat is $7$ cm long. Scale: $1 : 40$. Find the actual length. \answerBlank[3cm]
\answerExplain{$280$ cm}{Multiply: $7 \times 40 = 280$ cm. The actual boat is $280$ cm ($2.8$ m) long.}

\prob A model bridge is $15$ cm long. Scale: $1 : 300$. Find the actual length. \answerBlank[3cm]
\answerExplain{$4{,}500$ cm}{Multiply: $15 \times 300 = 4{,}500$ cm. The actual bridge is $4{,}500$ cm ($45$ m) long.}

\prob A model airplane is $20$ cm from wing tip to wing tip. Scale: $1 : 150$. Find the actual wingspan. \answerBlank[3cm]
\answerExplain{$3{,}000$ cm}{Multiply: $20 \times 150 = 3{,}000$ cm. The actual wingspan is $3{,}000$ cm ($30$ m).}
\end{multicols}
\end{practiceBox}

% ── Practice A: Actual to Model ──────────────────────────────────────────────
\begin{practiceBox}{Actual to Model (Divide by the Scale Factor)}

\prob A building is $40$ m tall. Scale: $1 : 500$. How tall is the model? \answerBlank[3cm]
\answerExplain{$8$ cm}{Convert to cm: $40$ m $= 4{,}000$ cm. Divide: $4{,}000 \div 500 = 8$ cm.}

\prob A road is $3$ km long. Scale: $1 : 10{,}000$. How long is the road on the model? \answerBlank[3cm]
\answerExplain{$30$ cm}{Convert: $3$ km $= 300{,}000$ cm. Divide: $300{,}000 \div 10{,}000 = 30$ cm.}

\prob A football field is $100$ m long. Scale: $1 : 250$. How long is the model? \answerBlank[3cm]
\answerExplain{$40$ cm}{Convert: $100$ m $= 10{,}000$ cm. Divide: $10{,}000 \div 250 = 40$ cm.}

\prob A ferry is $60$ m long. Scale: $1 : 200$. How long is the model? \answerBlank[3cm]
\answerExplain{$30$ cm}{Convert: $60$ m $= 6{,}000$ cm. Divide: $6{,}000 \div 200 = 30$ cm.}

\prob A monument is $25$ m tall. Scale: $1 : 500$. How tall is the model? \answerBlank[3cm]
\answerExplain{$5$ cm}{Convert: $25$ m $= 2{,}500$ cm. Divide: $2{,}500 \div 500 = 5$ cm.}
\end{practiceBox}

% ── Practice B: Finding the Scale Factor ─────────────────────────────────────
\begin{practiceBox}[funTeal]{What Is the Scale Factor?}

\prob A model car is $6$ cm long. The actual car is $4.8$ m long. Find the scale. \answerBlank[3cm]
\answerExplain{$1 : 80$}{Convert: $4.8$ m $= 480$ cm. Scale factor: $480 \div 6 = 80$. Scale is $1 : 80$.}

\prob A model airplane wing is $10$ cm. The real wing is $20$ m. Find the scale. \answerBlank[3cm]
\answerExplain{$1 : 200$}{Convert: $20$ m $= 2{,}000$ cm. Divide: $2{,}000 \div 10 = 200$. Scale is $1 : 200$.}

\prob A dollhouse room is $8$ cm wide. The actual room is $4$ m wide. Find the scale. \answerBlank[3cm]
\answerExplain{$1 : 50$}{Convert: $4$ m $= 400$ cm. Divide: $400 \div 8 = 50$. Scale is $1 : 50$.}

\prob A map distance of $2.5$ cm represents $50$ km. Find the scale. \answerBlank[3cm]
\answerExplain{$1 : 2{,}000{,}000$}{Convert: $50$ km $= 5{,}000{,}000$ cm. Divide: $5{,}000{,}000 \div 2.5 = 2{,}000{,}000$. Scale is $1 : 2{,}000{,}000$.}

\prob A model building is $12$ cm tall. The actual building is $36$ m tall. Find the scale. \answerBlank[3cm]
\answerExplain{$1 : 300$}{Convert: $36$ m $= 3{,}600$ cm. Divide: $3{,}600 \div 12 = 300$. Scale is $1 : 300$.}
\end{practiceBox}

% ── True or False ────────────────────────────────────────────────────────────
\begin{practiceBox}[funPurple]{True or False?}

\trueOrFalse{A scale of $1 : 100$ means the model is $100$ times \emph{larger} than the real object.}
\answerExplain{False}{A scale of $1 : 100$ means the model is $100$ times \emph{smaller} than the real object. Each cm on the model represents $100$ cm in real life.}

\trueOrFalse{If a model has scale $1 : 50$, a $10$ cm model length represents $5$ m in real life.}
\answerExplain{True}{$10 \times 50 = 500$ cm $= 5$ m. The actual length is $5$ m.}

\trueOrFalse{To go from actual size to model size, you multiply by the scale factor.}
\answerExplain{False}{To go from actual to model, you \emph{divide} by the scale factor. You multiply by the scale factor when going from model to actual.}
\end{practiceBox}

% ── Word Problems ────────────────────────────────────────────────────────────
\begin{practiceBox}[funBlue]{\faGlobe~Word Problems}

\wordProblem{An architect builds a $1 : 100$ scale model of a school. The model is $18$ cm long and $12$ cm wide. What are the actual dimensions of the school?}{m}
\answerExplain{$18$ m long by $12$ m wide}{Multiply each dimension by $100$: length $= 18 \times 100 = 1{,}800$ cm $= 18$ m; width $= 12 \times 100 = 1{,}200$ cm $= 12$ m.}

\wordProblem{A model of a skyscraper uses scale $1 : 400$. The actual building is $320$ m tall. How tall should the model be?}{cm}
\answerExplain{$80$ cm}{Convert: $320$ m $= 32{,}000$ cm. Divide: $32{,}000 \div 400 = 80$ cm.}

\wordProblem{A model train is $9$ cm long. The actual train engine is $18$ m long. Mia says the scale is $1 : 20$. Is she correct? Explain.}{}
\answerExplain{No, the scale is $1 : 200$}{Convert: $18$ m $= 1{,}800$ cm. Scale factor $= 1{,}800 \div 9 = 200$. The correct scale is $1 : 200$, not $1 : 20$. Mia forgot to convert meters to centimeters.}
\end{practiceBox}

% ── Challenge ────────────────────────────────────────────────────────────────
\begin{challengeBox}
\prob A $1 : 50$ model room is $10$ cm by $8$ cm. What is the actual area of the room in square meters? \answerBlank[3cm]
\answerExplain{$20$ m$^2$}{Actual length: $10 \times 50 = 500$ cm $= 5$ m. Actual width: $8 \times 50 = 400$ cm $= 4$ m. Area $= 5 \times 4 = 20$ m$^2$.}

\prob A model ship is built at two different scales. Scale A is $1 : 100$ and the model is $12$ cm long. Scale B makes the model $8$ cm long. What is Scale B? \answerBlank[3cm]
\answerExplain{$1 : 150$}{Actual length: $12 \times 100 = 1{,}200$ cm. Scale B factor: $1{,}200 \div 8 = 150$. Scale B is $1 : 150$.}
\end{challengeBox}

\encouragement{Great work with scale models! You can now calculate real-world sizes from any model!}
